% !TeX root = ../main.tex

\chapter{数据与变量}
\section{样本范围及可比性}
数据来自项目本地CSMAR合并表。股票收益保留市场类型1、4、16，估值截面沿用0、3、6开头的代码筛选。基础程序最后纳入2015年3月调仓，实际计收益至2015年6月；延展与四路径最后纳入2026年6月调仓，计收益至2026年9月。原始表覆盖终点与实际持有收益终点分别核查，不将未进入持仓计算的月份计入样本。

\begin{table}[htbp]
\centering
\caption{样本、字段与研究用途}
\renewcommand{\arraystretch}{1.18}
\small
\setlength{\tabcolsep}{4.5pt}
\begin{tabularx}{\textwidth}{@{}p{2.9cm}LL@{}}
\toprule
数据 & 主要字段或处理 & 用途 \\
\midrule
股票月度行情 & Mretwd；含现金红利再投资 & 组合月收益 \\
月末估值 & MarketValue；当月末总市值 & PE与规模 \\
净利润及EPS & B002000000；标准季度累计总利润 & 同比成长与TTM利润 \\
指数日行情 & 000001的Retindex；月内复利 & 上证综指基准 \\
冻结公告敏感性结果 & 2003-04至2024-12，261个月 & 附录信息时点诊断 \\
\bottomrule
\end{tabularx}
\end{table}

主比较分为两个层次。第一，147个月共同复现区间，基础复现、延展与四路径均使用相同月份。第二，282个月完整延展区间，以基础延展为同规则对照；基础复现不强行外推至2026年。年度表中2003年和2026年各只有9个月，其收益不当作完整自然年度收益。

\section{成长率、滚动利润与估值}
令NP为同一报告期的累计总净利润，成长率定义为：

\begin{equation}
g_{i,q}=\frac{NP_{i,q}}{NP_{i,q-4}}-1.
\label{eq:growth}
\end{equation}

去年同期利润非正时，成长率置为缺失，避免亏损扩大被机械识别为正增长。保留3、6、9、12月标准季末财报，不把1月1日的期初记录当作普通季度。非年报TTM利润由“上年全年利润$+$本期累计利润$-$去年同期累计利润”得到，年报本身直接使用全年利润。

\begin{equation}
PE_{i,t}=\frac{MV_{i,t}}{NP^{TTM}_{i,q}}.
\label{eq:pe}
\end{equation}

市值使用调仓当月值，利润与市值均有效且TTM非零时计算带符号PE。所有有限负PE参与升序排序并占用名次，最终入选要求PE为正。这严格沿用基础复现，并不等于先删除负PE、再在正PE内部取30\%。

\section{累计利润、TTM与同比的计算示例}
季度财报中的累计净利润包含年初至报告期末的利润，不能把一季度、半年度和前三季度累计数直接相加。举例而言，若上年全年净利润为100、今年半年度为70、上年同期半年度为50，则最近四个季度利润为100+70$-$50=120，累计利润同比为70/50$-$1=40\%。若同期总市值为2,400且单位相同，计算PE为20。这里的数字仅用于说明公式，不代表任一实际持仓；实际结果全部使用CSMAR原始字段。

同一条利润资料承担两种不同作用：同比衡量相对去年同期的成长，TTM为估值提供近一年盈利分母。二者不能互相替换。若上年同期为负，直接相除可能使扭亏、减亏或亏损扩大得到难以按通常正增长排序的数字，因此基础版将其同比置为缺失；缺失不表示增长为零，也不表示真实利润下降。历史“全为正”条件下，这条记录不能作为正成长通过。

市值除以总净利润与股价除以EPS具有相近的估值思想，但只有利润归属、股数定义及时间口径一致时才可严格互换。总净利润与归母净利润、期末总股本与EPS的加权平均股本均可能不同。本研究以基础程序口径为标准，不因另一字段提高收益就替换它；第8章的CSMAR归母利润、现成PE及行情表总市值对照，用于评估这些定义的敏感性。

\section{数据补齐解决了什么问题}
扩大时间覆盖能够增加可用月份和历史财报，避免延展结果因旧文件尾部缺失而停留在较早终点，但不自动解决每只股票的上市前财报覆盖、同比基数、财报修订和真实公告时间。对某一股票而言，“有20条报告记录”还不等于“有20个有效同比”：每一期同比还需要去年同期的利润基数，基数非正或缺失仍无法通过历史成长条件。

这也说明不能用“已补齐大部分时间区间”替代逐股诊断。本文同时检查输入版本、实际收益终点、财报连续性和同比有效数；原研报末期四只股票的缺口在第8章单独列明。当前证据只能说明本地资料按本规则会排除它们，不能推断这些公司当年真实经营不满足原作者条件。

\section{三维得分与财务信息边界}
质量指标为最近最多20条累计利润记录中净利润为正的比例，反映利润正值出现频率，不是ROE、现金流质量或单季利润率。动量为调仓月份之前12个月的累计收益，至少有8个月有效记录；规模为总市值，市值越小得分越高。三个指标在原有历史合格候选池内转换为方向一致的百分位后等权平均。

缺失质量按0、缺失动量按$-$1进入原有评分。四路径不改变此处理。财务指标直接按报告期映射到季度末，没有要求真实公告日已到达；使用最终会计值也未逐项恢复历史修订版本。因此，本文所有主表首先是既有基础口径下的回溯模拟。附录另以冻结的261个月公告敏感性结果说明这一限制。
